% Algebra 1 - ten-page chapter summary
% Source: Text & Tests 4 (New Edition), Chapter 1 "Algebra 1"
% Compile with: pdflatex algebra1_summary_10pp.tex
\documentclass[11pt]{article}
\usepackage[a4paper,top=1.9cm,bottom=1.9cm,left=2.1cm,right=2.1cm]{geometry}
\usepackage[T1]{fontenc}
\usepackage{amsmath,amssymb}
\usepackage{textcomp}
\usepackage{xcolor}
\usepackage{soul}
\usepackage{tcolorbox}
\usepackage{enumitem}
\usepackage{titlesec}
\usepackage{array,booktabs}
\usepackage{fancyhdr}
\usepackage[hidelinks]{hyperref}

\definecolor{lilac}{RGB}{228,214,248}      % light purple highlight
\definecolor{lilacdeep}{RGB}{104,58,160}   % headings
\definecolor{lilacline}{RGB}{176,150,214}  % rules and box frames
\definecolor{greyish}{RGB}{95,95,105}
\sethlcolor{lilac}

\newcommand{\key}[1]{\hl{#1}}
\newcommand{\keyin}[1]{\textbf{\color{lilacdeep}#1}}  % emphasis inside a purple box
% NB: maths inside \key{...} must be written literally as \mbox{$...$} (soul limitation)

\newtcolorbox{formulabox}[1][]{colback=lilac,colframe=lilacline,boxrule=0.7pt,
  arc=2pt,left=7pt,right=7pt,top=4pt,bottom=4pt,#1}
\newtcolorbox{methodbox}[1]{colback=white,colframe=lilacdeep,boxrule=0.7pt,arc=2pt,
  title={#1},coltitle=white,colbacktitle=lilacdeep,fonttitle=\bfseries\small,
  left=7pt,right=7pt,top=3pt,bottom=3pt}

\titleformat{\section}{\color{lilacdeep}\large\bfseries}{\thesection}{0.6em}{}
\titleformat{\subsection}{\color{lilacdeep}\normalsize\bfseries}{\thesubsection}{0.6em}{}
\titlespacing*{\section}{0pt}{1.0em}{0.35em}
\titlespacing*{\subsection}{0pt}{0.7em}{0.25em}

\pagestyle{fancy}\fancyhf{}
\renewcommand{\headrulewidth}{0pt}
\fancyfoot[L]{\footnotesize\color{greyish}Algebra 1 --- Chapter Summary (10 pages)}
\fancyfoot[R]{\footnotesize\color{greyish}\thepage}

\setlength{\parindent}{0pt}
\setlength{\parskip}{0.35em}

\newcommand{\wex}[1]{\textbf{\color{lilacdeep}Worked:} #1}

\begin{document}

{\color{lilacdeep}\Large\bfseries Algebra 1 --- Chapter Summary}\\[2pt]
{\color{greyish}\small \textit{Text \& Tests 4} (New Edition), Chapter 1 \quad$\cdot$\quad Leaving Certificate Higher Level \quad$\cdot$\quad Sections 1.1--1.10}\\[2pt]
{\color{greyish}\small Theory, methods and worked examples. Exercises are not reproduced here.}\\[-6pt]

\textcolor{lilacline}{\rule{\textwidth}{1pt}}

\key{Anything highlighted in light purple is a result or a rule you should know by heart.}
Boxes in purple are the formulae themselves; boxes with a purple banner are step-by-step methods.

\begin{formulabox}
\textbf{What the chapter covers}\\[3pt]
\begin{tabular}{@{}p{1.35cm}p{6.0cm}p{1.35cm}p{6.4cm}@{}}
\textbf{1.1} & Polynomial expressions: degree, expanding, dividing & \textbf{1.6} & Identities: equating coefficients, partial fractions\\
\textbf{1.2} & Polynomial functions: $f(x)$ notation, evaluating & \textbf{1.7} & Manipulating formulae: changing the subject\\
\textbf{1.3} & Factorising: five techniques & \textbf{1.8} & Patterns: first and second differences\\
\textbf{1.4} & Algebraic fractions & \textbf{1.9} & Solving equations: roots\\
\textbf{1.5} & Binomial expansions & \textbf{1.10} & Simultaneous equations in 2 and 3 unknowns\\
\end{tabular}\\[3pt]
The thread running through it: \keyin{every technique here is either a way of writing an expression as a product (factorising) or a way of turning one statement into an equation you can solve.}
\end{formulabox}

\subsection*{The chapter's key words}

\begin{tabular}{@{}p{2.6cm}p{13.2cm}@{}}
\textbf{expression} & a collection of terms, e.g.\ $5x^{3}-3x^{2}+4x-6$; it has a value once $x$ is known\\
\textbf{equation} & two expressions set equal, true only for particular values of the variable\\
\textbf{identity} & an equation true for \emph{every} value of the variable\\
\textbf{variable} & the letter whose value may change; the \textbf{independent} one is the one you choose\\
\textbf{degree} & the highest power appearing: linear (1), quadratic (2), cubic (3)\\
\textbf{expanding} & multiplying out brackets; \textbf{factorising} is the reverse\\
\textbf{factors} & expressions that multiply to give the original, leaving no remainder\\
\textbf{in terms of} & written using only the stated letters, e.g.\ ``$r$ in terms of $A$''\\
\textbf{simultaneous} & equations that must hold at the same time, for the same values\\
\end{tabular}

\section{Polynomial expressions (1.1)}

A \textbf{polynomial expression} is a sum of algebraic terms whose powers are \key{positive whole numbers}. For example $5x^{3}-3x^{2}+4x-6$ has four \textbf{terms}: $5x^{3}$, $-3x^{2}$, $4x$ and $-6$.

\begin{tabular}{@{}p{3.3cm}p{12.5cm}@{}}
\textbf{degree} & the highest power of $x$: degree 1 = \textbf{linear}, 2 = \textbf{quadratic}, 3 = \textbf{cubic}\\
\textbf{coefficient} & the number in front of a power: in $4x^{3}-2x^{2}+5x-6$ the coefficient of $x^{3}$ is 4, of $x^{2}$ is $-2$, of $x$ is 5\\
\textbf{constant term} & the term with no $x$ (here $-6$); also called the term independent of $x$\\
\textbf{binomial / trinomial} & an expression of two / three terms\\
\end{tabular}

Polynomials are written in order of decreasing powers ($5x^{3}-3x^{2}+4x-6$) or increasing powers ($9+3x-4x^{2}$).

\key{An expression is \emph{not} a polynomial if a variable appears in a denominator or under a root} --- e.g.\ $3x^{2}-\frac{4}{x}+x^{\frac32}$ fails twice, because $\frac4x=4x^{-1}$ has a negative power and $x^{3/2}$ a fractional one.

\subsection*{Adding, subtracting and multiplying}

Add and subtract by collecting like terms; multiply using the distributive law $a(b+c)=ab+ac$, taking every term of the first bracket by every term of the second. Multiplying out is called \textbf{expanding}.

\wex{$(x-5)(2x^{2}-3x+6)=x(2x^{2}-3x+6)-5(2x^{2}-3x+6)=2x^{3}-13x^{2}+21x-30$.}

\begin{formulabox}
\[(x+a)^{2}=x^{2}+2ax+a^{2}\qquad (x-a)^{2}=x^{2}-2ax+a^{2}\qquad (x-a)(x+a)=x^{2}-a^{2}\]
\end{formulabox}

The first two are \textbf{perfect squares}; the third is the \textbf{difference of two squares}, the single most useful factorising pattern in the chapter.

\textbf{Perfect-square problems.} To make $ax^{2}+bx+c$ a perfect square, match it against $(\sqrt a\,x\pm\sqrt c)^{2}$: the middle term must be twice the product of the two roots, i.e.\ \key{\mbox{$b=\pm 2\sqrt{ac}$}}.

\wex{$25x^{2}+px+16$ with $p>0$. Since $25x^{2}=(5x)^{2}$ and $16=(4)^{2}$, we need $(5x+4)^{2}$, so $p=2(5)(4)=40$. (Taking $16=(-4)^{2}$ gives $p=-40$, rejected because $p>0$.)}

\subsection*{Dividing}

Three situations arise:

\begin{enumerate}[topsep=2pt,itemsep=1pt,leftmargin=*]
\item \textbf{The denominator divides every term.} Split the fraction up:
$\dfrac{6x^{3}-8x^{2}y+4xy^{2}+2x^{2}}{2x}=3x^{2}-4xy+2y^{2}+x$.
\item \textbf{The denominator is a factor of the numerator.} Factorise and cancel:
$\dfrac{6x^{2}+5xy+y^{2}}{2x+y}=\dfrac{(3x+y)(2x+y)}{2x+y}=3x+y$.
\item \textbf{Otherwise, use long division.}
\end{enumerate}

\begin{methodbox}{Long division of polynomials}
\begin{enumerate}[topsep=1pt,itemsep=0pt,leftmargin=*]
\item Divide the first term of the divisor into the first term of what is left --- that gives the next term of the answer.
\item Multiply the whole divisor by that term and subtract.
\item Repeat with the new leading term until the remainder has lower degree than the divisor.
\end{enumerate}
\key{Write the dividend in decreasing powers and leave a gap for any missing power} (e.g.\ a cubic with no $x^{2}$ term), or the columns will not line up.
\end{methodbox}

\wex{$\dfrac{2x^{3}-9x^{2}+10x-3}{x-3}=2x^{2}-3x+1$, set out in full:}

\vspace{-0.4em}
\begin{center}
\begin{tabular}{@{}r@{\ }l@{\qquad}l@{}}
& \underline{$2x^{2}-3x+1$} & \\
$x-3\,\big)$ & $2x^{3}-9x^{2}+10x-3$ & \small\color{greyish}$2x^{3}\div x=2x^{2}$\\
& \underline{$2x^{3}-6x^{2}$}\hspace{2.55cm} & \small\color{greyish}multiply $2x^{2}$ by $(x-3)$, subtract\\
& $\ \ \ \ \ -3x^{2}+10x-3$ & \small\color{greyish}$-3x^{2}\div x=-3x$\\
& $\ \ \ \ \ \underline{-3x^{2}+9x}$\hspace{1.25cm} & \small\color{greyish}multiply, subtract\\
& $\ \ \ \ \ \ \ \ \ \ \ \ \ \ \ \ x-3$ & \small\color{greyish}$x\div x=1$\\
& $\ \ \ \ \ \ \ \ \ \ \ \ \ \ \ \ \underline{x-3}$ & \small\color{greyish}remainder $0$\\
\end{tabular}
\end{center}
\vspace{-0.3em}

A cubic with a missing power needs a gap: $(2x^{3}-11x+6)\div(2x^{2}+4x-3)=x-2$ is set out as $2x^{3}\ \underline{\hspace{6mm}}\ -11x+6$, with space left where the $x^{2}$ term would be.

The second of these also shows the link to factors: since the division is exact,
\[2x^{3}-11x+6=(x-2)(2x^{2}+4x-3),\]
so long division is a way of \emph{finding} factors, not just of dividing.

\section{Polynomial functions (1.2)}

If a rectangle has length $x$ and width $x-5$, its area depends on $x$; we write $A(x)=x(x-5)=x^{2}-5x$. Here $x$ is the \textbf{independent variable} and $A(x)$ the \textbf{dependent variable}.

\key{\mbox{$A(x)$} is one single symbol --- it does not mean \mbox{$A$} multiplied by \mbox{$x$}.}

\textbf{Evaluating} means substituting and simplifying. With $p(x)=2x^{2}-5x+6$:
\[p(1)=2-5+6=3,\qquad p(-3)=18+15+6=39,\qquad p(t)=2t^{2}-5t+6,\qquad p(t^{2})=2t^{4}-5t^{2}+6.\]

\key{You may substitute a whole expression, not just a number} --- $f(3a)$, $f(t-2)$ and $f\!\left(\frac a4\right)$ are all handled the same way, with brackets kept around the substituted expression.

\textbf{Combining functions.} Add or subtract by collecting like terms, then state the degree. With $f(x)=3x^{3}-4x^{2}-3x+4$ and $g(x)=5x^{3}+14x^{2}+7x-2$:
\[2f(x)-g(x)=x^{3}-22x^{2}-13x+10\ \ (\text{degree }3),\qquad f(x)+2g(x)=13x^{3}+24x^{2}+11x\ \ (\text{degree }3).\]

\wex{A rectangle of length $(2x+3)$ has area $A(x)=2x^{2}+7x+6$. The width is found by dividing: $\dfrac{2x^{2}+7x+6}{2x+3}=\dfrac{(2x+3)(x+2)}{2x+3}=x+2$, so the perimeter is $P(x)=2(2x+3)+2(x+2)=6x+10$.}

\wex{The daily cost of producing $x$ litres of paint is $C(x)=0.001x^{2}+0.1x+5$ (degree 2). Then $C(100)=10+10+5=\texteuro 25$ and $C(400)=160+40+5=\texteuro 205$: doubling and redoubling the output does not double the cost, because the model is not linear.}

\textbf{Context restricts the values.} A length must be positive, so if the width is $(x-5)$ then $x-5>0\Rightarrow x>5$; if the length is $(2x+3)$ then $x>-1.5$. \key{Always check what the variable is allowed to be when the question is about a real shape or a real quantity.}

\textbf{More than one variable.} $V(r,h)=\pi r^{2}h$ says the volume depends on both $r$ and $h$; its degree is 2, the highest power of \emph{either} variable.

\wex{An open box has dimensions $x+3$, $x+1$ and height $x$. Sides: $2x(x+3)+2x(x+1)=4x^{2}+8x$; base: $(x+3)(x+1)=x^{2}+4x+3$; so $S(x)=5x^{2}+12x+3$ and $S(5)=188\,\text{cm}^{2}$.}

\section{Factorising (1.3)}

\key{A factor divides into a polynomial leaving no remainder.} If $(x^{2}-x-12)\div(x-4)=x+3$, then $(x-4)$ and $(x+3)$ are both factors of $x^{2}-x-12$. Five techniques cover the chapter.

\textbf{1. Highest common factor.} $3x^{2}-9xy=3x(x-3y)$; $2a^{2}b-4ab^{2}+12abc=2ab(a-2b+6c)$.

\textbf{2. Grouping} (usually four terms, in two pairs):
\[6x^{2}y+3xy^{2}-12x-6y=3xy(2x+y)-6(2x+y)=(2x+y)(3xy-6)=3(2x+y)(xy-2),\]
where the last step is the reminder that a factor you have just produced may itself factorise.

\textbf{3. Difference of two squares.} Write both parts as squares first:
\begin{formulabox}
\[c^{2}-d^{2}=(c-d)(c+d),\qquad x^{2}-9y^{2}=(x-3y)(x+3y),\]
\[x^{2}-8y^{2}=\left(x-\sqrt8 y\right)\left(x+\sqrt8 y\right),\qquad x-9=\left(\sqrt x-3\right)\left(\sqrt x+3\right)\]
\end{formulabox}

\wex{$x^{4}-y^{4}=(x^{2}-y^{2})(x^{2}+y^{2})=(x-y)(x+y)(x^{2}+y^{2})$ --- one difference of squares can hide another.\\
$12x^{2}-75y^{2}=3(4x^{2}-25y^{2})=3(2x-5y)(2x+5y)$ --- take the common factor out first.}

\textbf{4. Quadratic trinomials $ax^{2}+bx+c$.} Either guess the factor pairs (trial and error), or, when the numbers are awkward or irrational, use the formula.

\begin{formulabox}
\[\text{If } ax^{2}+bx+c=0 \text{ then } x=\frac{-b\pm\sqrt{b^{2}-4ac}}{2a}.\qquad
\text{A root } x=k \text{ gives the factor } (x-k);\ \ x=\tfrac pq \text{ gives } (qx-p).\]
\end{formulabox}

\wex{$3x^{2}-17x+20$: the formula gives $x=4$ or $x=\frac53$, so the factors are $(x-4)$ and $(3x-5)$.\\
$x^{2}-2\sqrt2x-6$: $x=\dfrac{2\sqrt2\pm\sqrt{32}}{2}=\sqrt2\pm2\sqrt2=3\sqrt2$ or $-\sqrt2$, giving $\left(x-3\sqrt2\right)\left(x+\sqrt2\right)$.}

\textbf{5. Sum and difference of two cubes.} Obtained by long division of $x^{3}-y^{3}$ by $(x-y)$.

\begin{formulabox}
\[x^{3}-y^{3}=(x-y)(x^{2}+xy+y^{2})\qquad\qquad x^{3}+y^{3}=(x+y)(x^{2}-xy+y^{2})\]
\[(ax)^{3}-(by)^{3}=(ax-by)(a^{2}x^{2}+abxy+b^{2}y^{2})\qquad (ax)^{3}+(by)^{3}=(ax+by)(a^{2}x^{2}-abxy+b^{2}y^{2})\]
\end{formulabox}

Use them as templates: rewrite the expression as two cubes first, then substitute.

\wex{$a^{3}+8b^{3}=a^{3}+(2b)^{3}=(a+2b)(a^{2}-2ab+4b^{2})$.\\
$64c^{3}-125d^{3}=(4c)^{3}-(5d)^{3}=(4c-5d)\!\left[(4c)^{2}+(4c)(5d)+(5d)^{2}\right]=(4c-5d)(16c^{2}+20cd+25d^{2})$.}

\begin{methodbox}{Which technique to use}
\begin{enumerate}[topsep=1pt,itemsep=0pt,leftmargin=*]
\item Is there a common factor in every term? Take it out first --- always.
\item Two terms? Look for a difference of two squares, or a sum/difference of two cubes.
\item Three terms? It is a quadratic: trial and error, or the formula if the numbers are awkward.
\item Four terms? Group them in two pairs and look for a common bracket.
\item Then check each factor again: a factor may itself factorise.
\end{enumerate}
\key{Take out the highest common factor first, then keep going until nothing factorises further.}
\end{methodbox}

\section{Algebraic fractions (1.4)}

Algebraic fractions behave exactly like numerical ones.

\begin{formulabox}
\begin{enumerate}[topsep=1pt,itemsep=0pt,leftmargin=*]
\item A \textbf{common denominator} is needed to add or subtract.
\item A fraction can be reduced \textbf{only if the numerator and denominator share a common \emph{factor}}.
\item If a numerator or denominator itself contains added fractions, combine them into a single fraction first.
\item To divide, multiply by the reciprocal (the denominator turned upside down).
\end{enumerate}
\end{formulabox}

\key{Cancel factors, never terms.} In $\dfrac{6\times12}{3}$ you may cancel; in $\dfrac{6+12}{3}$ you may not cancel the 6 alone --- and the same is true with letters.

\textbf{Simplifying.} Factorise the top and the bottom fully, then cancel.
\[\frac{5ax}{15a+10a^{2}}=\frac{5a\,x}{5a(3+2a)}=\frac{x}{3+2a},\qquad
\frac{t^{2}+3t-4}{t^{2}-16}=\frac{(t+4)(t-1)}{(t-4)(t+4)}=\frac{t-1}{t-4},\qquad
\frac{x^{2}-9y^{2}}{3x+9y}=\frac{x-3y}{3}.\]

\textbf{Adding with factorised denominators.} Factorise \emph{before} choosing the common denominator, so the denominator is as small as possible.

\wex{$\dfrac{x-4}{x^{2}-x-2}-\dfrac{x-3}{x^{2}-4}
=\dfrac{x-4}{(x-2)(x+1)}-\dfrac{x-3}{(x-2)(x+2)}
=\dfrac{(x-4)(x+2)-(x-3)(x+1)}{(x-2)(x+1)(x+2)}=\dfrac{-5}{(x^{2}-4)(x+1)}$.}

\textbf{Complex fractions} (a fraction inside a fraction): turn the top into one fraction, turn the bottom into one fraction, then divide by inverting.

\wex{$\dfrac{\,y-\frac{x^{2}+y^{2}}{y}\,}{\ \frac1x-\frac1y\ }
=\dfrac{\frac{-x^{2}}{y}}{\frac{y-x}{xy}}
=\frac{-x^{2}}{y}\times\frac{xy}{y-x}=\frac{-x^{3}}{y-x}$.}

\section{Binomial expansions (1.5)}

Expanding $(a+b)^{n}$ by repeated multiplication is slow. The pattern of the expansions is:

\begin{enumerate}[topsep=2pt,itemsep=0pt,leftmargin=*]
\item powers of $a$ go \emph{down} by one from term to term, powers of $b$ go \emph{up} by one;
\item the powers in every term add to $n$;
\item there are \key{\mbox{$n+1$} terms} in the expansion of $(a+b)^{n}$;
\item the coefficients are the rows of \textbf{Pascal's triangle}, each entry the sum of the two above it:
\item those coefficients are the combinations $\binom nr={}^nC_r$, available on the calculator's \textsf{nCr} key.
\end{enumerate}

\vspace{-0.5em}
\begin{center}
\begin{tabular}{@{}r@{\quad}c@{\qquad}l@{}}
row 0 & $1$ & \small\color{greyish}$\binom00$\\
row 1 & $1\quad1$ & \small\color{greyish}$\binom10\ \binom11$\\
row 2 & $1\quad2\quad1$ & \small\color{greyish}$\binom20\ \binom21\ \binom22$\\
row 3 & $1\quad3\quad3\quad1$ & \small\color{greyish}coefficients of $(a+b)^{3}$\\
row 4 & $1\quad4\quad6\quad4\quad1$ & \small\color{greyish}coefficients of $(a+b)^{4}$\\
row 5 & $1\quad5\quad10\quad10\quad5\quad1$ & \small\color{greyish}coefficients of $(a+b)^{5}$\\
\end{tabular}
\end{center}

\begin{formulabox}
\textbf{Binomial Theorem} (page 20 of \textit{Formulae and Tables}):
\[(x+y)^{n}=\binom n0x^{n}+\binom n1x^{n-1}y+\binom n2x^{n-2}y^{2}+\cdots+\binom nrx^{n-r}y^{r}+\cdots+\binom nny^{n}\]
\textbf{General term:} $\displaystyle \binom nr x^{n-r}y^{r}$, with $r=0,1,2,\dots$
\end{formulabox}

Useful coefficient facts: $\binom n0=\binom nn=1$, $\binom n1=n$, and \key{\mbox{$\binom nr=\binom n{n-r}$}} (so $\binom 64=\binom 62$).

\key{Because the expansion starts at \mbox{$r=0$}, the \mbox{$(r+1)$}th term uses \mbox{$r$}: the 3rd term has \mbox{$r=2$}, the 4th term has \mbox{$r=3$}.} This off-by-one is the commonest slip in the whole section.

\wex{$(p+q)^{6}=p^{6}+6p^{5}q+15p^{4}q^{2}+20p^{3}q^{3}+15p^{2}q^{4}+6pq^{5}+q^{6}$.}

\wex{$(2x-3y)^{4}$: take $x\to 2x$, $y\to-3y$, $n=4$:
\[\binom40(2x)^{4}+\binom41(2x)^{3}(-3y)+\binom42(2x)^{2}(-3y)^{2}+\binom43(2x)(-3y)^{3}+\binom44(-3y)^{4}
=16x^{4}-96x^{3}y+216x^{2}y^{2}-216xy^{3}+81y^{4}.\]}

\key{Substitute the \emph{whole} term, sign and coefficient included, inside brackets} --- $(-3y)^{2}=+9y^{2}$, not $-9y^{2}$. When the second term is negative the signs alternate.

\wex{First three terms of $(1-5y)^{8}$: with $x=1$, $y=-5y$, $n=8$,
$1+\binom81(-5y)+\binom82(-5y)^{2}+\cdots=1-40y+700y^{2}+\cdots$}

\wex{Fourth term of $(3a+b)^{7}$: $r=3$, so $\binom73(3a)^{4}b^{3}=35\cdot81a^{4}b^{3}=2835a^{4}b^{3}$.}

\wex{Sixth term of $(2x+y)^{10}$: the 6th term means $r=5$, giving $\binom{10}{5}(2x)^{5}y^{5}=252\cdot32x^{5}y^{5}=8064x^{5}y^{5}$.}

To find one particular \emph{coefficient} rather than one particular term, do not expand everything: write the general term, decide which value of $r$ produces the power you want, and evaluate only that term.

\wex{The number of terms in $(p+2q)^{6}$ is $6+1=7$, so the middle term is the 4th, with $r=3$: $\binom63p^{3}(2q)^{3}=20\cdot8\,p^{3}q^{3}=160p^{3}q^{3}$.}

\textit{Note on the printed text:} in the worked example for $(1-5y)^{8}$ the book's third line prints $(1-3y)^{8}$; that is a typographical slip --- the arithmetic that follows ($-40y$, $700y^{2}$) is the expansion of $(1-5y)^{8}$.

\section{Algebraic identities (1.6)}

An \textbf{identity} is an equation that is true \key{for all values} of the variable, not just for one or two.

\begin{formulabox}
If $ax^{3}+bx^{2}+cx+d=px^{3}+qx^{2}+rx+s$ for all values of $x$, then
\[a=p,\qquad b=q,\qquad c=r,\qquad d=s.\]
\textbf{In an identity, the coefficients of like powers are equal} (and a power missing from one side has coefficient 0).
\end{formulabox}

So if $ax^{3}+bx^{2}+cx+d=qx^{2}+s$ for all $x$, then $a=c=0$, $b=q$ and $d=s$.

\wex{$(2x+a)^{2}=4x^{2}+12x+b$ for all $x$. Expanding: $4x^{2}+4ax+a^{2}=4x^{2}+12x+b$, so $4a=12\Rightarrow a=3$, and $a^{2}=b\Rightarrow b=9$.}

\wex{$3t^{2}x-3px+c-2t^{3}=0$ for all $x$. Write both sides as polynomials in $x$: $(3t^{2}-3p)x+(c-2t^{3})=(0)x+(0)$. Hence $3t^{2}=3p\Rightarrow t=\sqrt p$, and $c=2t^{3}=2p^{3/2}$.}

\subsection*{Partial fractions}

A single fraction can be split back into simpler ones, e.g.\ $\dfrac{1}{(x+1)(x-2)}=\dfrac{A}{x+1}+\dfrac{B}{x-2}$. Multiplying up gives the identity $1=A(x-2)+B(x+1)$, true for all $x$. Two methods:

\begin{methodbox}{Finding $A$ and $B$}
\textbf{Method 1 --- substitute convenient values.} Choose values of $x$ that kill one bracket:
$x=2\Rightarrow 1=3B\Rightarrow B=\frac13$; \ $x=-1\Rightarrow1=-3A\Rightarrow A=-\frac13$.\\[2pt]
\textbf{Method 2 --- equate coefficients.} $1=x(A+B)-2A+B$, so $A+B=0$ and $-2A+B=1$; solving simultaneously, $A=-\frac13$, $B=\frac13$.
\end{methodbox}
\[\therefore\quad \frac{1}{(x+1)(x-2)}=\frac{-1}{3(x+1)}+\frac{1}{3(x-2)}\]

An identity with three unknowns is handled the same way, and Method 1 is much the quicker.

\wex{$5x+3=Ax(x+3)+Bx(x-1)+C(x-1)(x+3)$ for all $x$. Put $x=0$: $3=-3C\Rightarrow C=-1$. Put $x=1$: $8=4A\Rightarrow A=2$. Put $x=-3$: $-12=12B\Rightarrow B=-1$.}

\subsection*{Identities and factors}

\key{If a divisor is a \emph{factor}, the division leaves no remainder --- so every coefficient of the remainder must be zero.} That turns a factor statement into equations connecting the unknown coefficients. Equivalently (and usually faster), write the missing factor with unknowns and expand:

\wex{$(x-t)^{2}$ is a factor of $x^{3}+3px+c$. Dividing by $x^{2}-2tx+t^{2}$ leaves the remainder $(3p+3t^{2})x+(c-2t^{3})$. Setting it identically to zero: $3p+3t^{2}=0\Rightarrow p=-t^{2}$, and $c-2t^{3}=0\Rightarrow c=2t^{3}$.}

\wex{If $x^{2}-ax+b$ is a factor of $x^{3}+2ax^{2}+4bx+c$, let the other factor be $(x+k)$. Expanding $(x+k)(x^{2}-ax+b)=x^{3}+(k-a)x^{2}+(b-ak)x+bk$ and comparing coefficients gives $k=3a$, $b=-a^{2}$ and $c=3ab$.}

\wex{$2x-\sqrt3$ is a factor of $4x^{2}-2(1+\sqrt3)x+\sqrt3$. Let the second factor be $(ax+b)$. Then $2ax^{2}+(2b-\sqrt3a)x-\sqrt3b$ must match, so $2a=4\Rightarrow a=2$; $2b-2\sqrt3=-2-2\sqrt3\Rightarrow b=-1$; the constants confirm it. The second factor is $2x-1$.}

\section{Manipulating formulae (1.7)}

Making $r$ the \textbf{subject} of $A=\pi r^{2}$ means writing $r$ in terms of the other letters: $r=\sqrt{A/\pi}$. The steps are always the same.

\begin{methodbox}{Changing the subject}
\begin{enumerate}[topsep=1pt,itemsep=0pt,leftmargin=*]
\item Clear roots (square both sides) and fractions (multiply by the denominators).
\item Expand any brackets.
\item Gather every term containing the required variable on one side, everything else on the other.
\item \key{If the variable appears more than once, factorise it out} --- this is the step people forget.
\item Divide by the bracket that multiplies it.
\end{enumerate}
\end{methodbox}

\wex{$v^{2}=u^{2}+2as\Rightarrow 2as=v^{2}-u^{2}\Rightarrow a=\dfrac{v^{2}-u^{2}}{2s}$.}

\wex{$\sqrt{\dfrac{x+y}{x-y}}=\dfrac12$. Squaring: $\dfrac{x+y}{x-y}=\dfrac14\Rightarrow4(x+y)=x-y\Rightarrow5y=-3x\Rightarrow y=\dfrac{-3x}{5}$; at $x=5$, $y=-3$.}

\wex{$x=\dfrac{t+4}{3t+1}$. Then $x(3t+1)=t+4\Rightarrow3tx-t=4-x\Rightarrow t(3x-1)=4-x\Rightarrow t=\dfrac{4-x}{3x-1}$.}

\wex{The pendulum formula $T=2\pi\sqrt{\dfrac lg}$. Divide by $2\pi$, then square: $\dfrac{T^{2}}{4\pi^{2}}=\dfrac lg$, so $l=\dfrac{gT^{2}}{4\pi^{2}}$. With $T=3$\,s and $g=10$\,m\,s$^{-2}$, $l=\dfrac{10(9)}{4\pi^{2}}\approx2.3$\,m.}

\key{Check a rearrangement by putting numbers into both versions} --- if the original and the rearranged form disagree, the slip is usually a sign lost when terms were moved across.

A formula can also be rewritten by \textbf{substituting} one relationship into another: for a cone with $\tan\theta=\frac rh$ we get $r=h\tan\theta$, so
\[V=\tfrac13\pi r^{2}h=\tfrac{\pi}{3}(h\tan\theta)^{2}h=\tfrac{\pi}{3}h^{3}\tan^{2}\theta .\]

\section{Algebraic patterns (1.8)}

\textbf{Linear patterns} ($0,3,6,9,\dots$ or $1,4,7,10,\dots$) go up by a constant amount.

\begin{formulabox}
For $f(x)=mx+c$ with $x=0,1,2,\dots$: \quad
\textbf{1st difference} $=m$ (the slope), \quad \textbf{starting value} $f(0)=c$ (the $y$-intercept).
\end{formulabox}

\textbf{Quadratic patterns} ($0,1,4,9,16,\dots$) do not. Their first differences ($1,3,5,7,\dots$) are not constant, but the \textbf{second differences} are.

\begin{formulabox}
For $f(x)=ax^{2}+bx+c$: \quad \keyin{the second difference is \mbox{$2a$}}. \quad
$f(x)=x^{2}+b$ is $\cup$-shaped, $f(x)=b-x^{2}$ is $\cap$-shaped, both symmetric about $x=0$.
\end{formulabox}

\begin{methodbox}{Finding the rule for a pattern}
\begin{enumerate}[topsep=1pt,itemsep=0pt,leftmargin=*]
\item Write out the 1st differences. Constant $\Rightarrow$ linear; if not, take the 2nd differences. Constant $\Rightarrow$ quadratic.
\item Quadratic: $c=f(0)$, the first term of the pattern; $a=\tfrac12(\text{2nd difference})$; then get $b$ from $f(1)=a+b+c$.
\item Check the rule against a term you have not used.
\end{enumerate}
\end{methodbox}

\wex{$4,7,16,31,52,79$.}

\vspace{-0.5em}
\begin{center}
\begin{tabular}{@{}l|cccccc@{}}
\textbf{pattern} & 4 & 7 & 16 & 31 & 52 & 79\\
\textbf{1st difference} & & 3 & 9 & 15 & 21 & 27\\
\textbf{2nd difference} & & & 6 & 6 & 6 & 6\\
\end{tabular}
\end{center}
\vspace{-0.3em}

The second difference is constant, so the pattern is quadratic: $2a=6\Rightarrow a=3$; $c=f(0)=4$; and $f(1)=3+b+4=7\Rightarrow b=0$. Hence $f(x)=3x^{2}+4$, which does give $f(2)=16$ and $f(3)=31$.

\wex{Matchstick patterns needing $4,10,18,28,\dots$ sticks: 2nd difference $=2\Rightarrow a=1$; $c=f(0)=4$; $f(1)=1+b+4=10\Rightarrow b=5$; so $f(x)=x^{2}+5x+4$. \key{The pattern starts at \mbox{$x=0$}, so the 10th pattern is \mbox{$f(9)=81+45+4=130$} matchsticks} --- not $f(10)$.}

\section{Solving equations (1.9)}

Solving an equation means finding the values of the variable that satisfy it. The values that make $y=0$ are the \textbf{roots}; \key{graphically, the roots are where the curve crosses the \mbox{$x$}-axis.} $4x-12=0$ has the single root $x=3$; $x^{2}-5x+6=0$ has two, $x=2$ and $x=3$.

Solving $2x+1=2$ instead of $2x+1=0$ is the same idea: it asks where the line $y=2x+1$ meets the line $y=2$. A linear equation always has exactly one root.

So there are three equivalent questions, and it is worth seeing them as one: \emph{solve $f(x)=0$}, \emph{find the roots of $f$}, and \emph{find where the graph of $f$ cuts the $x$-axis}. A linear function crosses once; a quadratic may cross twice, once, or not at all.

\begin{methodbox}{Linear equations containing fractions}
Find the lowest common denominator, multiply every term by it, expand, gather the variable on one side and the numbers on the other, then divide.
\end{methodbox}

\wex{$\dfrac{2t-3}{5}+\dfrac{1}{20}=\dfrac{t-1}{4}$. Multiply by 20: $4(2t-3)+1=5(t-1)\Rightarrow8t-12+1=5t-5\Rightarrow3t=6\Rightarrow t=2$.}

\section{Simultaneous equations (1.10)}

Two equations in $x$ and $y$ are either \textbf{parallel} (no common point) or meet at exactly one point $(x_1,y_1)$ that satisfies both \emph{simultaneously}. There are three ways to find it: substitution, elimination, or reading it off a graph.

\begin{methodbox}{Two unknowns}
\textbf{Substitution:} rearrange the easier equation to get one variable in terms of the other, substitute into the second equation, solve, then back-substitute.\\[2pt]
\textbf{Elimination:} multiply one or both equations so that one variable has matching coefficients, then add or subtract to remove it.\\[2pt]
\key{Always check the answer in \emph{both} original equations.}
\end{methodbox}

The graphical picture matters: solving a pair of equations is finding the point the two lines share. Parallel lines (same slope, different intercept) share no point and the algebra shows it --- both variables vanish and a false statement such as $0=7$ is left. If instead the two equations are multiples of one another, every point of the line is a solution.

\wex{(substitution) $3x-y=1$, $x-2y=-8$. From the second, $x=-8+2y$; then $3(-8+2y)-y=1\Rightarrow5y=25\Rightarrow y=5$, so $x=2$ and $(x,y)=(2,5)$.}

\wex{(elimination) $2x-5y=9$ (A), $3x+2y=4$ (B). $3\text{A}: 6x-15y=27$; $2\text{B}: 6x+4y=8$; subtracting, $-19y=19\Rightarrow y=-1$, then $2x-5(-1)=9\Rightarrow x=2$, so $(2,-1)$.}

\textbf{Three unknowns.} An equation such as $x+y+z=6$ is a \textbf{plane}. Three planes that are not parallel meet in a single point $(x,y,z)$.

\begin{methodbox}{Three unknowns}
\begin{enumerate}[topsep=1pt,itemsep=0pt,leftmargin=*]
\item Choose the unknown that is easiest to isolate.
\item Eliminate it twice, taking the equations \emph{two at a time}, to get two equations in two unknowns.
\item Solve those two as usual, then substitute back into one of the originals for the third unknown.
\item Check the solution in all three equations.
\end{enumerate}
\end{methodbox}

\wex{A: $x+y+z=6$; B: $2x+y-z=1$; C: $4x-3y+2z=4$. A${}+{}$B gives D: $3x+2y=7$; $2$B${}+{}$C gives E: $8x-y=6$. Then D${}+2$E: $19x=19\Rightarrow x=1$, so $y=2$ from D, and $z=3$ from A: $(x,y,z)=(1,2,3)$.}

Substitution is the method that carries forward: it is how the intersection of a line and a curve is found later in the course, so it is worth being fluent in it even where elimination is quicker.

\textbf{Problems in context.} The algebra is the easy part; the marks are in the setting up.

\begin{methodbox}{Turning words into equations}
State clearly what each letter stands for, write one equation for each piece of information given, \key{keep the units consistent} (e.g.\ \texteuro 32 becomes 3200 cent when the other terms are in cent), solve, and answer the question that was asked.
\end{methodbox}

\wex{240 people attended, tickets at \texteuro 31 and \texteuro 16, takings \texteuro 5595. With $x$ the number of \texteuro 16 tickets and $y$ the number of \texteuro 31 tickets: $x+y=240$ and $16x+31y=5595$. Eliminating $x$ gives $y=117$ and $x=123$.}

\wex{50c, 20c and 10c coins worth \texteuro 32 in total, with $2y+z=3x$ and $4x+z=6y$. In standard form: $50x+20y+10z=3200$, $3x-2y-z=0$, $4x-6y+z=0$. Eliminating $z$ leads to $x=40$, $y=35$, $z=50$ --- forty 50c, thirty-five 20c and fifty 10c coins.}

\subsection*{Where this chapter leads}

Nothing here is self-contained. The factorising and the quadratic formula are the tools for solving quadratic equations and for sketching curves; the binomial coefficients $\binom nr$ reappear as combinations in probability; first and second differences are taken up again in the chapter on sequences and series; and substitution into simultaneous equations becomes the standard way of finding where a line meets a curve. \key{Treat this chapter as the toolkit rather than as a topic --- almost every later chapter assumes it.}

\section*{\color{lilacdeep}Quick reference}

\begin{formulabox}
\[(x\pm a)^{2}=x^{2}\pm2ax+a^{2}\qquad (x-a)(x+a)=x^{2}-a^{2}\qquad x=\frac{-b\pm\sqrt{b^{2}-4ac}}{2a}\]
\[x^{3}\pm y^{3}=(x\pm y)(x^{2}\mp xy+y^{2})\qquad
(x+y)^{n}=\sum_{r=0}^{n}\binom nrx^{n-r}y^{r}\qquad \text{2nd difference}=2a\]
\end{formulabox}

\textbf{Revision checklist} --- can you, without looking anything up:
\begin{enumerate}[topsep=2pt,itemsep=0pt,leftmargin=*,label=$\square$]
\item state the degree, coefficients and constant term of a polynomial, and say why an expression is not one;
\item expand a product of two or three brackets, and recognise a perfect square and a difference of two squares;
\item divide one polynomial by another by long division, leaving gaps for missing powers, and read off the factors;
\item factorise by HCF, grouping, difference of two squares, quadratic trial or formula, and sum/difference of cubes;
\item add, subtract and simplify algebraic fractions, including complex ones;
\item write out any binomial expansion and pick out a named term or coefficient using $\binom nr x^{n-r}y^{r}$;
\item find unknown coefficients by equating like powers, and split a fraction into partial fractions;
\item change the subject of a formula when the subject appears twice;
\item decide from first and second differences whether a pattern is linear or quadratic, and find its rule;
\item solve simultaneous equations in two and three unknowns, and set them up from a worded problem?
\end{enumerate}

\textbf{Ten mistakes to avoid}

\begin{tabular}{@{}p{7.6cm}p{4.0cm}p{4.0cm}@{}}
\toprule
& \textbf{\color{lilacdeep}wrong} & \textbf{\color{lilacdeep}right}\\
\midrule
Calling an expression with a negative or fractional power a polynomial & $3x^{2}-\frac4x+x^{3/2}$ & polynomial powers must be $0,1,2,3,\dots$\\
Losing the gap for a missing power in a long division & $2x^{3}-11x+6$ & $2x^{3}\ \underline{\hspace{5mm}}\ -11x+6$\\
Cancelling a term instead of a factor & $\frac{6+12}{3}=6+4$ & $\frac{6+12}{3}=\frac{18}{3}=6$\\
Stopping halfway when factorising & $(x^{2}-y^{2})(x^{2}+y^{2})$ & $(x-y)(x+y)(x^{2}+y^{2})$\\
Missing the common factor first & $12x^{2}-75y^{2}$ ``won't factorise'' & $3(2x-5y)(2x+5y)$\\
Off-by-one in the binomial general term & 4th term uses $r=4$ & 4th term uses $r=3$\\
Dropping the bracket around a negative term & $-3y^{2}$ & $(-3y)^{2}=9y^{2}$\\
Not factorising out the subject when it appears twice & $3tx-t=4-x$, then stuck & $t(3x-1)=4-x$\\
Indexing a pattern from 1 when it starts at $x=0$ & 10th term $=f(10)$ & 10th term $=f(9)$\\
Ignoring what the variable is allowed to be & $x$ any real number & a width needs $x-5>0$\\
\bottomrule
\end{tabular}

\subsection*{What the \textit{Formulae and Tables} booklet gives you --- and what it does not}

The booklet supplies the \textbf{Binomial Theorem} and the notation $\binom nr={}^nC_r$ (page 20) and the \textbf{quadratic formula}. It does \emph{not} supply the expansions of $(x\pm a)^{2}$, the difference of two squares, the sum and difference of two cubes, or the fact that a constant second difference equals $2a$. \key{Those four have to be in your head.} Nor does it supply any method: the marks in this chapter go to the setting out --- the long-division layout, the line ``for all values of $x$, so the coefficients are equal'', the statement of what each letter stands for in a worded problem.

\subsection*{Symbols used in this chapter}

\begin{tabular}{@{}p{1.5cm}p{6.2cm}p{1.5cm}p{6.2cm}@{}}
$\Rightarrow$ & implies, ``and so'' & $\therefore$ & therefore\\
$\in \mathbb{N}$ & is a natural number $0,1,2,\dots$ & $\in \mathbb{R}$ & is a real number\\
$\binom nr$ & the number of ways of choosing $r$ objects from $n$ & $f(x)$ & the value of the function $f$ at $x$\\
\end{tabular}

\end{document}
